% ============================================================
%  Chapter 9 --- Buffer Management
% ============================================================
\chapter{Buffer Management}
\label{ch:buffer-management}

\section{The Buffer Pool Concept}

Disk I/O is orders of magnitude slower than memory access.  The \emph{buffer pool} is an in-memory cache of disk pages that exploits \emph{temporal locality} --- recently accessed pages are likely to be accessed again~\cite{tanenbaum2014modern}.

\begin{conceptbox}[The Buffer Pool Analogy]
Think of the buffer pool as a bookshelf next to your desk.  You can only keep $N$ books on the shelf.  When you need a new book, you fetch it from the library (disk).  If the shelf is full, you must return one book --- ideally the one you are least likely to need again.
\end{conceptbox}

\section{Buffer Pool Architecture in \shiftdb{}}

\begin{figure}[H]
\centering
\begin{tikzpicture}[node distance=0.8cm]
  % Buffer frames
  \draw[thick, fill=shiftBg] (0,0) rectangle (12,4.5);
  \node[font=\small\bfseries] at (6,4.2) {Buffer Pool (up to 1024 frames)};

  \foreach \i/\tbl/\pg/\dirty in {
    0/customers/0/clean,
    1/orders/2/dirty,
    2/products/0/clean,
    3/customers/1/dirty,
    4/orders/0/clean} {
    \pgfmathsetmacro{\x}{2.4*\i + 1.2}
    \node[slot, minimum width=2cm, minimum height=1.5cm] at (\x, 2) {
      Frame \i\\[2pt]
      \code{\tbl:\pg}\\[1pt]
      \textit{\dirty}
    };
  }

  % LRU list
  \node[font=\footnotesize\bfseries, anchor=west] at (0.3, 0.4) {LRU list (MRU $\to$ LRU):};
  \node[font=\scriptsize, anchor=west] at (5.5, 0.4) {orders:2 $\to$ customers:1 $\to$ products:0 $\to$ customers:0 $\to$ orders:0};
\end{tikzpicture}
\caption{Buffer pool with five frames showing table, page ID, and dirty status.}
\label{fig:buffer-pool}
\end{figure}

\subsection{BufferFrame Structure}

\begin{lstlisting}[style=cpp, caption={Buffer frame definition.}]
struct BufferFrame {
    Page page;            // 4 KB page data
    std::string tableName;
    uint32_t pageId = 0;
    bool dirty = false;   // modified since last write?
    int pinCount = 0;     // number of active users
};
\end{lstlisting}

\subsection{Key Operations}

\begin{table}[H]
\centering
\caption{Buffer pool operations.}
\begin{tabular}{l p{12cm}}
\toprule
\textbf{Method} & \textbf{Behaviour} \\
\midrule
\code{fetchPage}  & Return a pointer to the page.  On cache hit: bump LRU position, increment \code{pinCount}.  On miss: evict if pool is full, load from disk via \code{DiskManager::readPage}. \\
\code{markDirty}  & Set \code{dirty = true} for the given page.  Dirty pages must be flushed before eviction. \\
\code{unpinPage}  & Decrement \code{pinCount}.  A page with \code{pinCount = 0} is eligible for eviction. \\
\code{flushPage}  & Write a single dirty page to disk via \code{DiskManager::writePage}. \\
\code{flushAll}   & Iterate all frames and flush dirty pages (called on shutdown). \\
\code{evictTable} & Flush and remove all pages belonging to a dropped table. \\
\bottomrule
\end{tabular}
\end{table}

\section{LRU Eviction Policy}
\label{sec:lru}

\shiftdb{} uses the \textbf{Least Recently Used (LRU)} eviction policy, implemented with a \code{std::list} (doubly-linked list) tracking access order.

\begin{algorithm}[H]
\caption{LRU Eviction}
\label{alg:lru-eviction}
\KwInput{Buffer pool at capacity, need to load a new page}
\KwOutput{One page evicted}
Scan the LRU list from the \emph{back} (least recently used)\;
\ForEach{entry $e$ from LRU end}{
  \If{$e.\text{pinCount} = 0$}{
    \If{$e.\text{dirty}$}{
      Flush $e$ to disk\;
    }
    Remove $e$ from pool and LRU list\;
    \Return{}\;
  }
}
\tcp{All pages are pinned --- cannot evict (pool exhausted)}
\end{algorithm}

\section{Page Replacement Policies}

\begin{table}[H]
\centering
\caption{Comparison of page replacement policies.}
\begin{tabular}{l c p{8cm}}
\toprule
\textbf{Policy} & \textbf{Complexity} & \textbf{Trade-offs} \\
\midrule
LRU              & $\bigO{1}$ amortised & Good for temporal locality; vulnerable to sequential scans that flush useful pages \\
LRU-K            & $\bigO{\log n}$      & Tracks last $K$ access times; better resistance to scan pollution \\
Clock            & $\bigO{1}$ amortised & Approximation of LRU with lower overhead; used in PostgreSQL \\
2Q               & $\bigO{1}$           & Two queues (hot/cold); prevents one-time scans from polluting cache \\
ARC              & $\bigO{1}$           & Self-tuning between recency and frequency; used in ZFS \\
\bottomrule
\end{tabular}
\end{table}

\begin{deepdivebox}[Sequential Scan Problem]
If a query does \code{SELECT * FROM big\_table}, it touches every page exactly once.  With pure LRU, these pages push out frequently-used pages from other tables.  This is called the \emph{sequential flooding} problem.

\textbf{Solutions:}
\begin{itemize}[nosep]
  \item PostgreSQL uses a \emph{ring buffer} for large sequential scans --- only a few pages are used cyclically.
  \item LRU-K and 2Q give recently-accessed pages a ``probation'' period before promoting them to the hot list.
\end{itemize}
\end{deepdivebox}

\section{Thread Safety}

All buffer pool operations are protected by a \code{std::mutex}:
\begin{lstlisting}[style=cpp, caption={Mutex-protected fetchPage (excerpt).}]
Page* fetchPage(const std::string& tableName, uint32_t pageId) {
    std::lock_guard<std::mutex> lock(mutex_);
    // ... cache lookup and eviction logic
}
\end{lstlisting}

This coarse-grained lock is simple but limits concurrent access.  Production systems use finer-grained locking:
\begin{itemize}[nosep]
  \item \textbf{Per-frame latches} --- each \code{BufferFrame} has its own \code{std::shared\_mutex}.
  \item \textbf{Hash-table partitioning} --- divide the page map into shards, each with its own lock.
\end{itemize}

\section{Performance Metrics}

\shiftdb{} tracks buffer pool statistics:

\begin{lstlisting}[style=cpp, caption={Hit-rate calculation.}]
double getHitRate() const {
    uint64_t total = hits_ + misses_;
    return total > 0 ? static_cast<double>(hits_) / total : 0.0;
}
\end{lstlisting}

These metrics are exposed via the \code{GET /api/stats} endpoint and visualised in the frontend's Performance Monitor.

\begin{interviewbox}[Interview Corner]
\textbf{Q:} What is a buffer pool and why is it critical for DBMS performance?\\[4pt]
\textbf{A:} A buffer pool caches disk pages in memory to avoid slow disk reads.  A well-tuned buffer pool can serve 95\%+ of page requests from memory, making queries orders of magnitude faster.

\vspace{6pt}
\textbf{Q:} Explain the LRU page replacement policy.  What is its weakness?\\[4pt]
\textbf{A:} LRU evicts the page that was accessed least recently.  Its weakness is \emph{sequential flooding}: a full table scan touches every page once, evicting all useful cached pages.  Solutions include LRU-K, 2Q, or a dedicated scan buffer.

\vspace{6pt}
\textbf{Q:} What is a dirty page?  When is it flushed?\\[4pt]
\textbf{A:} A dirty page has been modified in memory but not yet written to disk.  It is flushed (1) upon eviction, (2) during a checkpoint, or (3) at shutdown.  The buffer pool must flush dirty pages before discarding them to avoid data loss.

\vspace{6pt}
\textbf{Q:} What is pin count and why does it matter for eviction?\\[4pt]
\textbf{A:} Pin count tracks how many active operations are using a page.  A page with \code{pinCount > 0} cannot be evicted --- doing so would pull the rug from under an active query.  Pages must be unpinned when the caller is done.
\end{interviewbox}
